% TOP_PROBLEM: {"id": 30007471, "title": "Two-fifths-power space simulation for one writable tape", "category_id": 15, "category": {"id": 15, "name": "computer_science", "display_name": "Computer Science", "description": "Computational complexity, algorithms, and theoretical CS.", "slug": "computer-science", "order_index": 15, "created_at": "2026-07-31T15:26:25.671Z"}, "status": "open"}
% FIRST_POSTED: 2026-10-06
% ATTEMPT_STATUS: solved
% OPENAI_CLAIM: {"schema_version":1,"source_repo":"https://github.com/openai/math","source_commit":"adc7f1241b42e322a6451854ab7e4b4c146bf78a","release_date":"2026-10-06","match":"full","resolution":"proved","independently_reviewed":false,"second_pass":"agreed","families":[{"family":"137","title":"One-tape time simulation in two-fifths-power space","subject":"Theoretical computer science","manuscripts":[{"title":"Simulating One-Tape Time in Two-Fifths-Power Space","pdf_path":"preprints/Simulating-One-Tape-Time-in-Two-Fifths-Power-Space-September-25-2026/article.pdf","readme_path":"preprints/Simulating-One-Tape-Time-in-Two-Fifths-Power-Space-September-25-2026/README.md","date":"2026-09-25","theorem_ref":"thm:main (source/sections/introduction.tex:64)","statement_sources":[{"path":"preprints/Simulating-One-Tape-Time-in-Two-Fifths-Power-Space-September-25-2026/build/source/sections/introduction.tex","line":64,"label":"thm:main"}]}],"lean":null,"reasoning_trace":null}]}
\documentclass[11pt]{article}
\usepackage[margin=25mm]{geometry}
\usepackage{fontspec}
\setmainfont{Latin Modern Roman}
\usepackage{amsmath,amssymb,mathtools}
\usepackage{unicode-math}
\setmathfont{Latin Modern Math}
\usepackage{enumitem}
\usepackage{xurl,hyperref,bookmark}
\hypersetup{colorlinks=true,urlcolor=blue}
\setcounter{secnumdepth}{0}
\setlength{\parindent}{0pt}
\setlength{\parskip}{5pt}
\setlength{\emergencystretch}{3em}
\begin{document}
\section*{Two-fifths-power space simulation for one writable tape}

\subsection{OpenAI's claim}
OpenAI presents the following claims relevant to this question, with manuscript numbers following paper order in Sources. In manuscript 1 (family 137; 2026-09-25), the asserted result is that a fixed one-writable-tape machine run for a supplied cap $T$ can be simulated in $O(T^{2/5}\operatorname{polylog}(T+2))$ work bits, with the stated initial-symbol accessor and unrestricted time.

\subsection{Mathematical statement of the claimed result}
\paragraph{Manuscript 1.}
Fix a deterministic finite-control machine with one integer-indexed writable tape/head and finitely many read-only heads on the input, with their number fixed, all moving by at most one cell per step from fixed origins. Initial symbol functions $\sigma_h(x,q)$ are fixed independently of time caps. The access premise requires one fixed deterministic accessor and constants $C_0,C_1$ such that, for every tag $h$ and $|q|\le T$, it returns $\sigma_h(x,q)$ using at most $C_0\log^{C_1}(T+2)$ work bits. Given binary cap $T\ge2$ satisfying that premise, the claimed simulator determines the control state and whether the machine halts within $T$ steps, using
\[
 O(T^{2/5}\log^C(T+2))
\]
work bits for a fixed $C$, excluding read-only input storage. There is no asserted bound on simulation time. Outputting the final contents of the writable tape is not required. If the access premise holds for every cap and the original machine halts after $t$ steps, the analogous outcome uses $O((t+2)^{2/5}\log^C(t+2))$ space without a supplied running time.

\subsection{Outline of the argument}
For manuscript 1 (family 137; 2026-09-25), block visits and epoch checkpoints create dependencies with different storage weights. An interpolation evaluator borrows and restores registers from arbitrary contents, so only designated dependency edges save large indices; balancing three scales gives the two-fifths exponent.

\subsection{Scope relative to this problem}
The result supplies the conclusion in this problem with its displayed hypotheses and quantifiers. The broader statements and applications in the manuscript are kept separate from the particular question above. The supporting result is manuscript 1 (family 137; 2026-09-25).

\subsection{Formal verification}
No Lean formalization is listed at this snapshot for manuscript 1 (family 137; 2026-09-25).

This site has not executed the Lean code. Described statement coverage is not an independent proof check.

\subsection{Status caveat}
The manuscripts were produced by an internal OpenAI model. OpenAI cautions that results without formalization can contain mistakes. This site has not independently checked the proofs; the label records the scope of OpenAI's claim.

\subsection{Sources}
\begin{itemize}
\item OpenAI, \emph{Simulating One-Tape Time in Two-Fifths-Power Space}, 2026-09-25; citation key \texttt{OAI:Simulating-One-Tape-Time-in-Two-Fifths-Power-Space-September-25-2026}. \href{https://github.com/openai/math/blob/adc7f1241b42e322a6451854ab7e4b4c146bf78a/preprints/Simulating-One-Tape-Time-in-Two-Fifths-Power-Space-September-25-2026/article.pdf}{PDF (GitHub)}; \href{https://github.com/openai/math/raw/adc7f1241b42e322a6451854ab7e4b4c146bf78a/preprints/Simulating-One-Tape-Time-in-Two-Fifths-Power-Space-September-25-2026/article.pdf}{PDF (direct)}; \href{https://github.com/openai/math/blob/adc7f1241b42e322a6451854ab7e4b4c146bf78a/preprints/Simulating-One-Tape-Time-in-Two-Fifths-Power-Space-September-25-2026/README.md}{Manuscript page and citation}; \href{https://github.com/openai/math/blob/main/preprints/Simulating-One-Tape-Time-in-Two-Fifths-Power-Space-September-25-2026/article.pdf}{Latest version}; Statement source: \href{https://github.com/openai/math/blob/adc7f1241b42e322a6451854ab7e4b4c146bf78a/preprints/Simulating-One-Tape-Time-in-Two-Fifths-Power-Space-September-25-2026/build/source/sections/introduction.tex\#L64}{thm:main}.
\item \href{https://github.com/openai/math/blob/adc7f1241b42e322a6451854ab7e4b4c146bf78a/CONTENTS.md}{OpenAI catalogue entry}.
\end{itemize}
\end{document}
